\documentclass{article}
\usepackage[utf8]{inputenc}

\usepackage{amsmath, amssymb, amsthm, amsfonts}
\usepackage{subcaption, enumerate}
\usepackage{bbold}
\usepackage{polynom}
\usepackage{xfrac}
\usepackage{graphicx}
\usepackage{mathtools}
\usepackage{enumitem}

\usepackage{tikz}

\newcommand{\dab}
{
	\begin{tikzpicture}[scale=0.04]
		% head
		\draw[thick] (0, 0) circle (3cm);

		% bent arm
		\draw[thick, rotate around={15 : (0, 0)}] (-3, 0) -- (-3.7, -2) -- (4.1, -1.1);

		% extended arm
		\draw[thick, rotate=20] (3, 0) -- (7, 2);
	\end{tikzpicture}
}
\newcommand{\st}{\text{ s.t. }}

\usepackage{fullpage}

\title{Linear Transforms Day 3 \\ with Dr Nevard}
\author{Notes by Laura Gallo}
\date{September 22, 2020}

\begin{document}
\maketitle

\section{Proposition: If $T:V\rightarrow W$ is linear and bijective, $T^{-1}$ is linear.}
\begin{proof}
	Let $T:V\rightarrow W$, $v,w\in V$
	$$TT^{-1}(av+bw)=av+bw$$
	$$T(aT^{-1}v+bT^{-1}w)=aTT^{-1}v+bTT^{-1}w=av+bw$$
	By transitivity, $TT^{-1}(av+bw)=T(aT^{-1}v+bT^{-1}w)$ \\
	$\therefore T^{-1}(av+bw)=aT^{-1}v+bT^{-1}w$ \dab
\end{proof}

\section{$\mathcal{L}(V,W)$}
If $V,W$ are vector spaces, then we define $\mathcal{L}(V,W)=\{\text{all linear transforms } T:V\rightarrow W\}$.

\subsection{Examples}
\begin{itemize}[label=--]
	\item If $V$ is a vector space over the scalar field $\mathbb{F}$ (which is also a vector space), then $\mathcal{L}(V,\mathbb{F})$ is called the space of functionals.
	\item $\mathcal{L}(V)=\mathcal{L}(V,V)$ is an algebra because for all $S,T,R\in \mathcal{L}(V)$:
		\begin{itemize}[label=--]
			\item $cS$, $S+T$, and $ST$ are all defined and closed in $\mathcal{L}(V)$
			\item $R(ST)=(RS)T$
			\item $R(S+T)=RS+RT$, $(S+T)R=SR+TR$ (in general $ST\neq TS$)
		\end{itemize}
\end{itemize}

\section{Powers of Linear Transforms}
In $\mathcal{L}(V)$, $T^0=I$, $T^n=TT^{n-1}$. \\
With this, we can have:
\begin{itemize}[label=--]
	\item Polynomials of $T$: $\sum_{k=0}^n a_kT^k$
	\item Power series of $T$: $\sum_{n=0}^\infty a_nT^n$ (must be tested for convergence)
	\item Exponents: $e^T=\sum_{n=0}^\infty \frac{T^n}{n!}$ (always converges)
\end{itemize}

\section{Properties of Injective Linear Transforms}
If and only if $T:V\rightarrow W$ is an injective linear transform, the following hold:
\begin{itemize}[label=--]
	\item The nullity of $T$ is 0
		\begin{proof}
			Left to right follows trivially from injectivity. \\
			For right to left, suppose only $T(0)=0$
			$$Tv_1=Tv_2$$
			$$T(v_1-v_2)=0$$
			$$v_1-v_2=0$$
			$$v_1=v_2$$
			$\therefore T$ is injective. \dab
		\end{proof}
	\item $T$ preserves linear independence
		\begin{proof}
			Let $\{v_n\}\subset V$ be a linearly independent set.
			$$c_1Tv_1+c_2Tv_2+...+c_nTv_n=0$$
			$$T(c_1v_1+c_2v_2+...+c_nv_n)=0$$
			$$c_1v_1+c_2v_2+...+c_nv_n=0$$
			Because ${v_n}$ is linearly independent, $c_k=0$. \\
			$\therefore \{Tv_n\}$ is linearly independent. \dab
		\end{proof}
	\item $\text{dim } V=\text{dim }T(V)$ \\
		This follows from the above and the Big Ugly Lemma.
\end{itemize}

\section{Representing Linear Transforms}
For $T\in \mathcal{L}(V,W)$, if $\{v_1,v_2,...,v_n\}\subset V$ a basis, $Tv_k=w_k$, then $\exists!c_k\st v=c_1v_1+c_2v_2+...+c_nv_n$.
In general, given two bases $\{v_n\}\subset V,\{w_m\}\subset W$, we can describe the linear transform $T:V\rightarrow W$ in the following way:

$$Tv_1=a_{11}w_1+a{21}w_2+...+a_{m1}w_m$$
$$Tv_2=a_{12}w_1+a{22}w_2+...+a_{m2}w_m$$
$$\vdots$$
$$Tv_n=a_{1n}w_1+a{2n}w_2+...+a_{mn}w_m$$

We define the matrix of $T$ with respect to the bases $\{v_n\}$ and $\{w_m\}$ in the following way:
\begin{displaymath}
	\mathcal{M}(\{v_n\},\{w_m\}) = \left(\begin{matrix}
		a_{11} & a_{12} & \cdots & a_{1n} \\
		a_{21} & a_{22} & \cdots & a_{2n} \\
		\vdots & \vdots & \ddots & \vdots \\
		a_{m1} & a_{m2} & \cdots & a_{mn} \\
	\end{matrix}\right)
\end{displaymath}

To apply this transformation onto $v=c_1v_1+c_2v_2+...+c_nv_n$, we do the following:
$$Tv=T(c_1v_1+c_2v_2+...+c_nv_n)=cTv_1+cTv_2+...+cTv_n=\sum_{k=1}^nc_kTv_k$$
$$=\sum_{k=1}^nc_k\sum_{j=1}^ma_{jk}w_k=\sum_{j=1}^mw_k\sum_{k=1}^nc_ka_{jk}$$

\end{document}
